\section{Stable nonlinear dynamics}\label{sec:dynamics}

The certificate construction extends to a class of equality-constrained
problems with a feasible forward repair. Local affine strips contain the
nonlinear graphs. Backward adjoints choose separator slopes whose
first-order terms vanish on every coordinate changed by the repair.

\subsection{The constrained model and graph strips}

Consider states $\sigma_0,\ldots,\sigma_T$ and continuous controls
$u_0,\ldots,u_{T-1}$, where $T\ge1$, on the product box
\[
 X=\prod_{t=0}^T I_t\times\prod_{t=0}^{T-1}U_t.
\]
The intervals are compact and nondegenerate. The symbol $s_0$ retains its
earlier meaning as the largest interval width. The feasible set is
\begin{equation}\label{eq:dynamics-feasible}
 P=\{x\in X:q_t(x_{V_t})=0,\ 0\le t<T\},\qquad
 q_t(\sigma,u,\sigma')=\sigma'-\varphi_t(\sigma,u).
\end{equation}
There are no additional constraints. The decomposition is the path rooted
at bag zero, with
\[
 V_t=\{\sigma_t,u_t,\sigma_{t+1}\},\qquad
 S_t=\{\sigma_t\}\quad(1\le t<T).
\]
Assign $q_t$ and the bag cost $a_t$ to $V_t$, and put $F=\sum_ta_t$.
Write $q(x)=(q_t(x_{V_t}))_{t=0}^{T-1}$ for the residual vector.
Here $N=T$, $p=3$, and we use the uniform occurrence bound $k=2$ even
when $T=1$. The largest separator dimension is one when $T\ge2$ and zero
when $T=1$.

\begin{assumption}\label{ass:dynamics}
The maps $\varphi_t$ are continuously differentiable on neighborhoods of
their input rectangles, and the following bounds hold throughout those
rectangles:
\begin{equation}\label{eq:dynamics-promises}
 \varphi_t(I_t\times U_t)\subseteq I_{t+1},\qquad
 |\partial_\sigma\varphi_t|\le a<1,\qquad
 |\partial_u\varphi_t|\le b,\qquad
 \Lip(\nabla\varphi_t)\le H,
\end{equation}
where $a,b,H\ge0$. The costs and their continuous convex models satisfy
\eqref{eq:model-smooth} and \eqref{eq:model-contract} on the full bag
boxes. There are $x^*\in P$ and $g>0$ such that
\begin{equation}\label{eq:dynamics-growth}
 F(y)-F(x^*)\ge g\norm{y-x^*}^2\qquad(y\in P).
\end{equation}
At every feasible center the dependent-state objective derivatives satisfy
\begin{equation}\label{eq:dynamics-state-gradient}
 |\partial_{\sigma_j}F(c)|\le G
 \qquad(c\in P,\ 1\le j\le T).
\end{equation}
\end{assumption}

The growth condition concerns the true feasible set. Stability alone does
not imply it. A finite $G$ exists for each fixed compact instance; a
horizon-uniform conclusion requires a horizon-uniform bound. If every
bag-gradient component has magnitude at most $G_b$ on its full box, then
$G=2G_b$ is valid on all of $X$, including infeasible centers. The graph
curvature bound $H$ is used to construct the following valid local domains
and must be supplied correctly.

\begin{lemma}[Affine graph strips]\label{lem:dynamics-strip}
Let $B\subseteq X_{V_t}$ have width $w_B$, and let $\xi_B$ be the
midpoint of its projection onto the input coordinates. Define
\begin{align}\label{eq:dynamics-strip}
 \ell_{t,B}(\sigma,u)
 &=\varphi_t(\xi_B)+\nabla\varphi_t(\xi_B)^{\mathsf T}
                                  ((\sigma,u)-\xi_B),\\
 O_t(B)&=\{z\in B:
 |z_{\sigma_{t+1}}-\ell_{t,B}(z_{\sigma_t},z_{u_t})|
                                     \le Hw_B^2/4\}.\nonumber
\end{align}
The set $O_t(B)$ is a compact convex polytope, possibly empty. It contains
every true graph point in $B$, and every $z\in O_t(B)$ satisfies
\begin{equation}\label{eq:dynamics-strip-residual}
 |q_t(z)|\le K w_B^2,\qquad K=H/2.
\end{equation}
\end{lemma}
\begin{proof}
The integral Taylor remainder and the two input-coordinate widths give
\[
 |\varphi_t(\sigma,u)-\ell_{t,B}(\sigma,u)|
 \le\frac H2\norm{(\sigma,u)-\xi_B}^2\le\frac H4w_B^2.
\]
This proves graph containment. For a strip point, the triangle inequality
adds its strip deviation, at most $Hw_B^2/4$, to this Taylor error and
gives \eqref{eq:dynamics-strip-residual}. All constraints defining the
strip are affine.
\end{proof}

Strips on neighboring leaves may differ on their common faces. Each
contains the true graph points of its own leaf, which suffices for
certificate validity. If $H=0$, the strip is the exact affine graph and
may have smaller dimension.

For this constrained problem, replace the local domains of
\Cref{def:model-certificate} and \eqref{eq:model-intercept} by
\begin{equation}\label{eq:dynamics-local-domain}
 C_{t,B,D}=O_t(B)\cap\{z:z_{S_t}\in D\}.
\end{equation}
At the root omit the separator condition. The objective is still
$m_{t,B}$ plus affine separator terms; it is convex on this polytope.
The algorithm does not minimize on the nonlinear graph or construct a
convex model of a multiplier-adjusted objective.

Empty domains are represented by infeasibility flags. A flagged separator
cell asserts that no feasible subtree assignment has its separator in that
cell. A bag leaf is unusable when all intersecting cells of some child
are flagged. Otherwise its child intercept is the minimum over the
unflagged intersecting cells. An own cell is flagged when every incident
leaf is unusable or has empty local domain. Only finite intercepts are
stored for unflagged cells.

These assertions and the lower bounds are valid by bottom-up induction.
A true feasible assignment selects a local graph point in its leaf strip
and own cell. Its child separator value selects an intersecting cell that
cannot be flagged by the induction hypothesis. Therefore its leaf is
usable and its own local domain is nonempty. On each such assignment the
usual affine child-minorant and local inequalities apply, proving the
subtree lower bound and, at the root, $\LB\le F(x^*)$. Every true
trajectory induces an available consistent configuration, so a finite
root configuration exists. Continuity ensures attainment on nonempty
compact domains. The unfolding proof of \Cref{prop:model-dp} then applies
to finite choices with these domains and flags.

The exact local oracle returns either the minimum and an attaining point
on a nonempty domain \eqref{eq:dynamics-local-domain}, or a certificate
that the domain is empty. An empty return counts as an oracle call.

\subsection{Forward repair and adjoint cancellation}

Define $y=R(x)$ by preserving $\sigma_0$ and every control and recursively
setting
\begin{equation}\label{eq:dynamics-repair}
 y_{\sigma_{t+1}}=\varphi_t(y_{\sigma_t},x_{u_t}).
\end{equation}
The invariance promise in \eqref{eq:dynamics-promises} proves by induction
that $R:X\to P$ and that $R$ fixes every point in $P$.

\begin{lemma}[Stable repair]\label{lem:dynamics-repair}
For every $x\in X$,
\begin{equation}\label{eq:dynamics-repair-bound}
 \norm{R(x)-x}\le\frac{\norm{q(x)}}{1-a}.
\end{equation}
\end{lemma}
\begin{proof}
Set $e_t=R(x)_{\sigma_t}-x_{\sigma_t}$ and
$r_t=q_t(x_{V_t})$. Both state arguments lie in $I_t$, so
\[
 e_0=0,\qquad
 |e_{t+1}|\le a|e_t|+|r_t|.
\]
Extend the residual magnitudes by zero. The error-magnitude vector is
bounded coordinatewise by a sum of lag shifts with coefficients
$1,a,a^2,\ldots$. Each shift has Euclidean norm at most one. Taking norms
and using the geometric sum proves \eqref{eq:dynamics-repair-bound}.
\end{proof}

\begin{lemma}[Adjoints at an ambient center]\label{lem:dynamics-adjoint}
Let $c\in X$ satisfy $|\partial_{\sigma_j}F(c)|\le G$ for
$1\le j\le T$, whether or not $c$ is feasible. Set
\[
 d_j=\partial_{\sigma_j}F(c),\qquad
 A_t=\partial_\sigma\varphi_t(c_{\sigma_t},c_{u_t}),
\]
and define the backward recurrence
\begin{equation}\label{eq:dynamics-adjoint}
 \mu_{T-1}=-d_T,\qquad
 \mu_{t-1}=A_t\mu_t-d_t\quad(t=T-1,\ldots,1).
\end{equation}
For $a_t^c=a_t+\mu_tq_t$ and $G_c=\sum_ta_t^c$,
\begin{equation}\label{eq:dynamics-adjoint-cancel}
 |\mu_t|\le U:=\frac G{1-a},\qquad
 \partial_{\sigma_j}G_c(c)=0\quad(1\le j\le T),\qquad
 \nabla G_c(c)^{\mathsf T}(x-R(x))=0\quad(x\in X).
\end{equation}
Each $a_t^c$ has gradient-Lipschitz constant at most
$\overline M=M+UH$. Its subtree-gradient separator slopes are
\begin{equation}\label{eq:dynamics-slopes}
 \lambda_t(c)=\partial_{\sigma_t}a_t(c_{V_t})-\mu_t A_t
 \qquad(1\le t<T).
\end{equation}
\end{lemma}
\begin{proof}
The backward recurrence gives
$|\mu_t|\le G\sum_{r=0}^{T-1-t}a^r\le U$. Direct differentiation
gives $\partial_{\sigma_T}G_c(c)=d_T+\mu_{T-1}=0$ and, for
$1\le j<T$,
\[
 \partial_{\sigma_j}G_c(c)=d_j+\mu_{j-1}-A_j\mu_j=0.
\]
This uses no equation $q(c)=0$. Every repair displacement is supported
on $\sigma_1,\ldots,\sigma_T$, proving the final cancellation identity.
The gradient of $q_t$ is
$(-\partial_\sigma\varphi_t,-\partial_u\varphi_t,1)$, with Lipschitz
constant at most $H$; hence the adjusted gradient has constant at most
$M+UH$. Within the subtree of bag $t$, its separator state $\sigma_t$
occurs only in that bag, giving \eqref{eq:dynamics-slopes}.
\end{proof}

The adjusted gradient annihilates the entire coordinate subspace containing
all repair displacements. No linearity or differentiability of $R$ is
needed. This stronger cancellation also permits rational, infeasible
centers in the next section. The derivative bound is needed because graph
curvature magnifies multiplier terms; objective curvature alone does not
bound multipliers induced by large affine costs. The obstruction in
\Cref{prop:limitations-constraint} illustrates the failure of unmodified
objective-gradient slopes under constraints.

\subsection{Aggregate error and constructive bounds}

Use the following constants, with $N=T$, $p=3$, $k=2$:
\begin{equation}\label{eq:dynamics-constants}
 \begin{aligned}
 C_0&=k(k-1)p,& U&=G/(1-a),&\overline M&=M+UH,\\
 K&=H/2,&L_q&=\sqrt{1+a^2+b^2},&
 \rho&=\frac{L_q\sqrt{C_0}+Ks_0}{1-a},\\
 D&=(\sqrt{C_0}+\sqrt{k}\rho)^2,&
 B_0&=\overline M D/2+A/4+UK,&\eta&=g/20,\\
 C&=12B_0+6kD\overline M^2/\eta,&
 B&=\max\{p,2C/g\}.&&
 \end{aligned}
\end{equation}
Choose a dyadic $\theta\le1/2$ satisfying
\begin{equation}\label{eq:dynamics-grading}
 12D\theta^2\le1,\qquad
 \sqrt{12}\,k\sqrt D\,\overline M\theta\le\eta/4,\qquad
 12kB_0\theta^2\le\eta/4.
\end{equation}
Zero coefficients impose no restriction.

\begin{lemma}[Configuration error after repair]\label{lem:dynamics-error}
Build the shell partitions of \Cref{lem:regridding-shell} around a center
$c\in X$ satisfying the derivative bound of
\Cref{lem:dynamics-adjoint}, at scale $h$ and ratio $\theta$ satisfying
\eqref{eq:dynamics-grading}. Use slopes \eqref{eq:dynamics-slopes}.
For every finite configuration with consistent top-copy point $x$,
forward repair $y=R(x)$, and configuration value $\Phi$,
\begin{equation}\label{eq:dynamics-error-final}
 |F(y)-\Phi|\le\frac g{20}\norm{y-c}^2+CNh^2.
\end{equation}
\end{lemma}
\begin{proof}
Let $b_t=\width(B_t)$, $d_t=\width(D_t)$, and put
\[
 Q=\sum_t b_t^2+\sum_{t=1}^{T-1}d_t^2,\quad
 E_x=\sum_t\norm{z^t-x_{V_t}}^2,\quad
 E_y=\sum_t\norm{z^t-y_{V_t}}^2.
\]
The incidence-based argument of \Cref{lem:regridding-drift} gives
$E_x\le C_0Q$ independently of the strip constraints. Each residual
$q_t$ is $L_q$-Lipschitz on its full bag box. By
\eqref{eq:dynamics-strip-residual} and $b_t\le s_0$,
\[
 \norm{q(x)}\le L_q\sqrt{E_x}+K\sqrt{\sum_t b_t^4}
             \le(L_q\sqrt{C_0}+Ks_0)\sqrt Q.
\]
\Cref{lem:dynamics-repair} gives $\norm{y-x}\le\rho\sqrt Q$.
Stacking the bag vectors and using coordinate multiplicity yields
\begin{equation}\label{eq:dynamics-copy-repair}
 \sqrt{E_y}\le\sqrt{E_x}+\sqrt{k}\norm{y-x},\qquad E_y\le DQ.
\end{equation}
The term $Ks_0$ remains necessary even for one bag: a strip point may
require repair although it has no copy mismatch.

Write $\operatorname{err}_t=a_t(z^t)-m_{t,B_t}(z^t)$. Apply
\Cref{lem:model-telescope} to the adjusted gradients to obtain
\[
 \Phi=\sum_ta_t(z^t)-\sum_t\operatorname{err}_t
       -\sum_t\nabla a_t^c(c_{V_t})^{\mathsf T}(z^t-x_{V_t}).
\]
\Cref{lem:dynamics-adjoint} permits replacing $x$ by $y$ in the last
sum. Because $q_t(y_{V_t})=0$, rearrangement gives the exact identity
\begin{equation}\label{eq:dynamics-identity}
 \Phi=F(y)+\sum_tT_t-\sum_t\operatorname{err}_t
                              -\sum_t\mu_tq_t(z^t),
\end{equation}
where
\[
 T_t=a_t^c(z^t)-a_t^c(y_{V_t})
                -\nabla a_t^c(c_{V_t})^{\mathsf T}(z^t-y_{V_t}).
\]
The multiplier residual in \eqref{eq:dynamics-identity} accounts for
the relaxed graphs and has absolute value at most $UKQ$.
The integral gradient remainder bounds
\[
 |T_t|\le\overline M\norm{y_{V_t}-c_{V_t}}\norm{z^t-y_{V_t}}
                     +\frac{\overline M}{2}\norm{z^t-y_{V_t}}^2.
\]
Set $r=\norm{y-c}$. By Cauchy--Schwarz,
\eqref{eq:model-contract}, and \eqref{eq:dynamics-copy-repair},
\begin{equation}\label{eq:dynamics-error-q}
 |F(y)-\Phi|\le\overline M\sqrt{kD}\,r\sqrt Q+B_0Q.
\end{equation}

For a leaf containing $z^t$, the shell width bound gives
$b_t\le h+\theta\norm{y_{V_t}-c_{V_t}}
                +\theta\norm{z^t-y_{V_t}}$.
The same holds for its own cell with coordinates restricted to $S_t$.
Each coordinate occurs in at most $k$ bags and $k-1$ separators.
The bag-plus-separator copy-error sum is at most $2E_y$.
It need not equal $2E_y$, since repair may change private top-bag states.
Summing $(v_1+v_2+v_3)^2\le3\sum_i v_i^2$ therefore gives
\[
 Q\le6Nh^2+3(2k-1)\theta^2r^2+6\theta^2E_y.
\]
By \eqref{eq:dynamics-copy-repair} and the first grading condition,
absorption yields
\begin{equation}\label{eq:dynamics-aggregate-grading}
 Q\le12Nh^2+12k\theta^2r^2.
\end{equation}
Insert this bound into \eqref{eq:dynamics-error-q}, using
$\sqrt{v+w}\le\sqrt v+\sqrt w$. The result is
\[
 |F(y)-\Phi|\le\sqrt{12kD}\,\overline M\sqrt N\,hr
 +[\sqrt{12}\,k\sqrt D\,\overline M\theta+12kB_0\theta^2]r^2
 +12B_0Nh^2.
\]
Young's inequality bounds the mixed term by
$\eta r^2/2+6kD\overline M^2Nh^2/\eta$; the remaining grading
conditions bound the bracket by $\eta/2$. This proves
\eqref{eq:dynamics-error-final}.
\end{proof}

\begin{theorem}[Constructive certificates for stable dynamics]
\label{thm:dynamics}
Under \Cref{ass:dynamics}, fix a dyadic ratio satisfying
\eqref{eq:dynamics-grading}. Start with any $y_{-1}\in P$, obtained by
forward repair if needed, and initialize $\UB=F(y_{-1})$. At stage $j$,
build fresh shell partitions at
$c=y_{j-1}$ and $h_j=s_0 2^{-j}$, construct the strips, compute
\eqref{eq:dynamics-adjoint} and \eqref{eq:dynamics-slopes}, and run the
flagged convex dynamic program. Backtrack a minimizing configuration,
form its top-copy point $x_j$, and set $y_j=R(x_j)$. Update the incumbent
with $F(y_j)$ and stop when $\UB-\LB_j\le\eps$.

Every stopping output is valid. With the constants
\eqref{eq:dynamics-constants}, termination occurs by
\begin{equation}\label{eq:dynamics-stage}
 J=\max\{0,\ceil{\log_2(s_0\sqrt{gBN/\eps})}\}.
\end{equation}
For $m=4/\theta$ and largest actual separator dimension $q$, the counts
are bounded by
\begin{equation}\label{eq:dynamics-counts}
 \begin{aligned}
 \text{final leaves and cells}&\le2Nm^p(J+1),\\
 \text{total created boxes}&\le Nm^p(J+1)(J+2),\\
 \text{local convex calls}&\le
 N3^qm^p(J+1)(J+2)(2J+3)/6.
 \end{aligned}
\end{equation}
Every local problem has at most three variables and affine constraints.
\end{theorem}
\begin{proof}
The local-domain induction already proves $\LB_j\le F(x^*)$ and the
existence of a finite minimizing configuration. Put
$e_j=\norm{y_j-x^*}^2$. Feasible growth and
\Cref{lem:dynamics-error} give
\[
 ge_j\le F(y_j)-\LB_j
 \le\frac g{20}\norm{y_j-y_{j-1}}^2+CNh_j^2
 \le\frac g{10}(e_j+e_{j-1})+CNh_j^2.
\]
Hence $e_j\le e_{j-1}/9+(10C/(9g))Nh_j^2$.
Initially $e_{-1}\le n s_0^2\le pNh_0^2$.
Since $B\ge p$ and $B\ge2C/g$, the stage-zero estimate implies
$e_0\le BNh_0^2$. For $j\ge1$, the induction closes because
$h_{j-1}=2h_j$ and $4B/9+10C/(9g)\le B$. Thus
$e_j\le BNh_j^2$ throughout.

For $j\ge1$, $\norm{y_j-y_{j-1}}^2\le10BNh_j^2$; at stage zero
the bound $4BNh_0^2$ suffices. Therefore
\[
 0\le\UB-\LB_j\le F(y_j)-\LB_j
      \le(gB/2+C)Nh_j^2\le gBNh_j^2.
\]
This proves \eqref{eq:dynamics-stage}. The stopping rule remains valid
at ratios outside the contraction conditions, since each lower certificate
and feasible upper value remains valid independently of those conditions.

The shell counts and \Cref{prop:regridding-incidence} give
\eqref{eq:dynamics-counts}. Graph strips preserve the incidence lists,
and empty local domains still count as calls. Each strip requires one
value and gradient evaluation of its map; each stage also requires
$O(T)$ evaluations and arithmetic for objective state derivatives,
adjoints, and forward repair. These costs fit the displayed work bounds
when charged as exact oracle operations. Costs of evaluating arbitrary
lists of factors assigned to a bag must be charged separately.
\end{proof}

For uniformly bounded numerical constants
$a,b,H,G,M,A,g,s_0$, with $1-a$ and $g$ uniformly positive, the theorem
gives $O(T(1+\log_+(T/\eps)))$ final boxes and cells, the squared
logarithm for all created boxes, and the cubed logarithm for local oracle
calls, where $\log_+z=\max\{0,\log z\}$. These are exact-real oracle counts.
Nonlinear simulation may produce
long rational coordinates or transcendental states; the theorem does not
bound their expanded representation or the internal cost of convex
minimization. The next section supplies a finite-precision result for
fixed-degree rational polynomials with compressed trajectory output.

\begin{example}[Uniform nonlinear conditioning]\label{ex:dynamics}
Let every interval be $[-1,1]$ and set
\[
 \varphi_t(\sigma,u)=\sigma/4+u/2+\sigma^2/8,\qquad
 F=\sigma_0^2+\sum_{t=0}^{T-1}
           [\sigma_{t+1}^2+u_t^2-u_t^4/4].
\]
Assign the initial square to bag zero. The map image is
$[-5/8,7/8]$, and $a=b=1/2$, $H=1/4$, $M=2$, and $G=2$ are valid.
Since $u^2-u^4/4\ge3u^2/4$ on $[-1,1]$, the zero trajectory is the
unique minimizer and $g=3/4$ works even on the full box. Keep square
factors exact and, on a control interval $[l,r]$, replace its quartic cost
by $u^2-u^4/4-(u-l)(r-u)/2$. Its second derivative is
$3-3u^2\ge0$ and its model error is at most $(r-l)^2/8$, so $A=1/2$.
All these constants are independent of $T$.

The reduced objective is nonconvex: set the initial state and all earlier
controls to zero and vary only the final control $v$. Then
$\sigma_T=v/2$ and the second derivative of the restricted objective is
$5/2-3v^2$, which is negative at $v=15/16$. Keeping the state variables
retains three-variable bags and one-variable separators under nonlinear
dynamics.
\end{example}

The invariance promise ensures feasibility for every retained initial state
and control sequence in their intervals. A terminal target or additional
state restriction requires a separate repair argument. Enumerating discrete
controls can require exponential work and is outside the horizon counts
above.
